\newpage
\section{第26章：有界量化}

\begin{taplauthorsolutionentrybox}
\phantomsection\label{sol:author:26.2.3}
\textbf{26.2.3 解答}

Abadi、Cardelli 与 Viswanathan（1996）的对象编码中确实需要完全
\(F_{<:}\) 的量词规则；Abadi 与 Cardelli（1996）也讨论了这一例子。
\taplsolutionbacklink{ex:26.2.3}{返回练习26.2.3}
\end{taplauthorsolutionentrybox}

\begin{taplauthorsolutionentrybox}
\phantomsection\label{sol:author:26.3.4}
\textbf{26.3.4 解答}
\[
\begin{aligned}
\taplterm{spluszz}
={}&\lambda n:\tapltype{SZero}.\lambda m:\tapltype{SZero}.
\lambda X.\lambda S<:X.\lambda Z<:X.\lambda s:X\to S.\lambda z:Z.\\
&n[X][S][Z]s\,(m[X][S][Z]s\,z),\\
\taplterm{spluspn}
={}&\lambda n:\tapltype{SPos}.\lambda m:\tapltype{SNat}.
\lambda X.\lambda S<:X.\lambda Z<:X.\lambda s:X\to S.\lambda z:Z.\\
&n[X][S][X]s\,(m[X][S][Z]s\,z).
\end{aligned}
\]
两者类型分别为
\(\tapltype{SZero}\to\tapltype{SZero}\to\tapltype{SZero}\) 和
\(\tapltype{SPos}\to\tapltype{SNat}\to\tapltype{SPos}\)。
\taplsolutionbacklink{ex:26.3.4}{返回练习26.3.4}
\end{taplauthorsolutionentrybox}

\begin{taplauthorsolutionentrybox}
\phantomsection\label{sol:author:26.3.5}
\textbf{26.3.5 解答}
\[
\begin{aligned}
\tapltype{SBool}
 &=\forall X.\forall T<:X.\forall F<:X.T\to F\to X,\\
\tapltype{STrue}
 &=\forall X.\forall T<:X.\forall F<:X.T\to F\to T,\\
\tapltype{SFalse}
 &=\forall X.\forall T<:X.\forall F<:X.T\to F\to F.
\end{aligned}
\]
\begin{lstlisting}
tru = lambda X. lambda T<:X. lambda F<:X.
      lambda t:T. lambda f:F. t;
? tru : STrue

fls = lambda X. lambda T<:X. lambda F<:X.
      lambda t:T. lambda f:F. f;
? fls : SFalse
\end{lstlisting}
\[
\begin{aligned}
\taplterm{notft}
 &=\lambda b:\tapltype{SFalse}.\lambda X.\lambda T<:X.\lambda F<:X.
   \lambda t:T.\lambda f:F.\,b[X][F][T]f\,t,\\
\taplterm{nottf}
 &=\lambda b:\tapltype{STrue}.\lambda X.\lambda T<:X.\lambda F<:X.
   \lambda t:T.\lambda f:F.\,b[X][F][T]f\,t.
\end{aligned}
\]
\taplsolutionbacklink{ex:26.3.5}{返回练习26.3.5}
\end{taplauthorsolutionentrybox}

\begin{taplauthorsolutionentrybox}
\phantomsection\label{sol:author:26.4.3}
\textbf{26.4.3 解答}

赋型推导的 T-Abs、T-TAbs 情形需要把新增绑定移到适合应用归纳假设的位置；
子类型推导的 S-All 情形也需要这样做。类型替换相关证明中的量词情形同样依赖
排列。
\taplsolutionbacklink{ex:26.4.3}{返回练习26.4.3}
\end{taplauthorsolutionentrybox}

\begin{taplauthorsolutionentrybox}
\phantomsection\label{sol:author:26.4.5}
\textbf{26.4.5 解答}

第 1 项对子类型推导归纳。S-Refl、S-Top 立即成立；S-Trans、S-Arrow、
S-All 直接使用归纳假设。S-TVar 有两个子情形。若最后查询的变量 \(Y\neq X\)，
原上界绑定在收窄后的上下文中仍存在。若 \(Y=X\)，原结论使用 \(X<:Q\)；
收窄后先由 S-TVar 得 \(X<:P\)，再由弱化把假设
\(\Gamma\vdash P<:Q\) 搬入扩展上下文，最后用 S-Trans 得 \(X<:Q\)。

第 2 项对赋型推导归纳，并在 T-TApp 的子类型前提以及 T-Sub 情形使用第 1 项。
\taplsolutionbacklink{lem:26.4.5}{返回引理26.4.5}
\end{taplauthorsolutionentrybox}

\begin{taplauthorsolutionentrybox}
\phantomsection\label{sol:author:26.4.11}
\textbf{26.4.11 解答}

四项均对子类型推导归纳。以箭头情形为例：S-Refl、S-Top 直接给出结论；
S-TVar 和 S-All 不可能产生箭头左侧；S-Arrow 的两个子推导正是所需关系。
若末步为 S-Trans，先对左子推导使用归纳假设，得到中间类型为
\(\tapltype{Top}\) 或某个箭头。前者再用第 4 项可知目标也是
\(\tapltype{Top}\)；后者再对右子推导使用归纳假设，并用两次传递性连接参数
和结果位置的关系。全称类型情形同理。
\taplsolutionbacklink{ex:26.4.11}{返回练习26.4.11}
\end{taplauthorsolutionentrybox}

\begin{taplauthorsolutionentrybox}
\phantomsection\label{sol:author:26.5.1}
\textbf{26.5.1 解答}

完全版本允许两个表示类型的上界不同，并令上界协变：
\[
\inferrule*[right=S-Some（完全）]
 {\Gamma\vdash S_1<:T_1\\
  \Gamma,X<:S_1\vdash S_2<:T_2}
 {\Gamma\vdash
  \{\exists X<:S_1,S_2\}<:\{\exists X<:T_1,T_2\}}.
\]
\taplsolutionbacklink{ex:26.5.1}{返回练习26.5.1}
\end{taplauthorsolutionentrybox}

\begin{taplauthorsolutionentrybox}
\phantomsection\label{sol:author:26.5.2}
\textbf{26.5.2 解答}

没有子类型时恰有四种：
\begin{lstlisting}
{*Nat, {a=5,b=7}} as {Some X, {a:Nat,b:Nat}};
{*Nat, {a=5,b=7}} as {Some X, {a:X,b:Nat}};
{*Nat, {a=5,b=7}} as {Some X, {a:Nat,b:X}};
{*Nat, {a=5,b=7}} as {Some X, {a:X,b:X}};
\end{lstlisting}
加入子类型和有界量化后还会出现许多选择，例如：
\begin{lstlisting}
{*Nat, {a=5,b=7}} as {Some X, {a:Nat}};
{*Nat, {a=5,b=7}} as {Some X, {b:X}};
{*Nat, {a=5,b=7}} as {Some X, {a:Top,b:X}};
{*Nat, {a=5,b=7}} as {Some X, Top};
{*Nat, {a=5,b=7}} as {Some X<:Nat, {a:X,b:X}};
{*Nat, {a=5,b=7}} as {Some X<:Nat, {a:Top,b:X}};
\end{lstlisting}
\taplsolutionbacklink{ex:26.5.2}{返回练习26.5.2}
\end{taplauthorsolutionentrybox}

\begin{taplauthorsolutionentrybox}
\phantomsection\label{sol:author:26.5.3}
\textbf{26.5.3 解答}

一种办法是把重置计数器 ADT 嵌套在普通计数器 ADT 中：
\begin{lstlisting}
counterADT =
  {*Nat,
   {new = 1,
    get = lambda i:Nat. i,
    inc = lambda i:Nat. succ(i),
    rcADT =
      {*Nat,
       {new = 1,
        get = lambda i:Nat. i,
        inc = lambda i:Nat. succ(i),
        reset = lambda i:Nat. 1}}
      as {Some ResetCounter<:Nat,
          {new:ResetCounter,
           get:ResetCounter->Nat,
           inc:ResetCounter->ResetCounter,
           reset:ResetCounter->ResetCounter}}}}
  as {Some Counter,
      {new:Counter,
       get:Counter->Nat,
       inc:Counter->Counter,
       rcADT:
         {Some ResetCounter<:Counter,
          {new:ResetCounter,
           get:ResetCounter->Nat,
           inc:ResetCounter->ResetCounter,
           reset:ResetCounter->ResetCounter}}}};
? counterADT
  : {Some Counter,
     {new:Counter,get:Counter->Nat,inc:Counter->Counter,
      rcADT:
        {Some ResetCounter<:Counter,
         {new:ResetCounter,get:ResetCounter->Nat,
          inc:ResetCounter->ResetCounter,
          reset:ResetCounter->ResetCounter}}}}
\end{lstlisting}
依次打开这两个包后，余下程序的上下文会含有
\(\tapltype{Counter}<:\tapltype{Top}\) 与
\(\tapltype{ResetCounter}<:\tapltype{Counter}\)，以及两套操作记录：
\begin{lstlisting}
let {Counter,counter} = counterADT in
let {ResetCounter,resetCounter} = counter.rcADT in
counter.get
  (counter.inc
    (resetCounter.reset
      (resetCounter.inc resetCounter.new)));
? 2 : Nat
\end{lstlisting}
\taplsolutionbacklink{ex:26.5.3}{返回练习26.5.3}
\end{taplauthorsolutionentrybox}

\begin{taplauthorsolutionentrybox}
\phantomsection\label{sol:author:26.5.4}
\textbf{26.5.4 解答}

在\cref{sec:24.3}的编码中给隐藏类型的量词加上同一上界：
\[
\llbracket\{\exists X<:S,T\}\rrbracket
=\forall Y.(\forall X<:\llbracket S\rrbracket.
\llbracket T\rrbracket\to Y)\to Y.
\]
打包和开包翻译只需相应地把类型抽象写成
\(\lambda X<:\llbracket S\rrbracket\)。由有界全称类型的子类型规则逐层
推导，可得到 S-Some。
\taplsolutionbacklink{ex:26.5.4}{返回练习26.5.4}
\end{taplauthorsolutionentrybox}
